\section{Assessment.cls}\label{sec:ref-assessment}

The graded-document class: tests, finals, quizzes, homework, worksheets.
Article-based. Loads Exercises, ProblemMeta, Workspace, Typefaces (lining
figures), MathStuff and InstructorNotes; deliberately \emph{not} Titling --
an assessment opens with \cs{AssessmentTitle}, a different artefact from a
styled title. Every page carries a ``Page $N$ of $M$'' footer (two-pass via
\file{lastpage}) and no running head. The text block stays wide, near
article's stock margins: an exam is not continuous prose, so the room goes
to the student's writing, not to a reading measure. Workflow:
section~\ref{sec:assessments}.

\subsection*{Class options}

\begin{center}
\begin{tabular}{@{}l p{0.62\linewidth}@{}}
\texttt{noid} & no student identity collected (homework, worksheets, MC
  practice): drops the NAME/ID box and the deep scanner top margin.\\
\texttt{automatch} & Crowdmark auto-matching: keeps the deep top margin,
  replaces the NAME/ID box with a blank 9--18\,cm mid-cover band.\\
\texttt{tuesday}, \texttt{thursday} & framed day box beside the title, TUE
  left / THU right -- sort two section piles by sight.\\
\texttt{legalpaper} & legal stock; absolute margins, so the extra height is
  writing room.\\
\emph{(other)} & forwarded to \file{article} (\texttt{12pt},
  \texttt{twoside}, \dots).\\
\end{tabular}
\end{center}

The default (no options) is the fully-furnished identified exam --
forgetting an option over-furnishes homework harmlessly, never under-equips
an exam.

\subsection*{Commands}

\begin{center}
\begin{tabular}{@{}l p{0.55\linewidth}@{}}
\cs{assessmentsetup}\texttt{\{...\}} & keyval cover setter: \env{course},
  \env{term}, \env{name}, \env{date}, \env{instructions}; all optional,
  unset fields vanish cleanly.\\
\cs{AssessmentTitle} & emits the cover: NAME/ID box (unless suppressed),
  centred title, instructions; adds a \emph{SOLUTIONS} stamp in a solutions
  build.\\
\cs{BonusMarksOnTop} & bonus margin block as ``(N pts)'' over rule over
  ``bonus'' (the default).\\
\cs{BonusLabelOnTop} & the flip: ``BONUS'' over rule over ``(N pts)''.\\
\cs{BonusMarginLabel}, \cs{BonusMarginRule} & reword the bonus label /
  restyle the divider rule.\\
\cs{ExamTopMargin} & the deep-top length (3.7\,cm), single-sourced for
  geometry and the \texttt{[automatch]} band; rarely touched.\\
\end{tabular}
\end{center}

\subsection*{Presentation choices}

The class overrides the Exercises hooks (section~\ref{sec:hooks}) to the
exam look: questions render as \textbf{Q$N$} hung in the left margin
(\cs{ProblemLabelWord} is \texttt{Q}; \cs{ProblemNumberText} is the tight
``Q$N$'', no space), and every scored item drops its ``($N$ pts)'' mark
into the right margin rail (\cs{ProblemPointsFormat} via \cs{marginnote}) --
numbers left, marks right, so standalone questions and lettered group parts
line up in the same rails. Bonus items stack their mark and label in that
same right rail (\cs{ProblemBonusPointsFormat}). \texttt{number=} and
\texttt{desc=} import keys are ignored: an exam owns its numbering and does
not annotate its questions.

One geometry fact worth knowing rather than rediscovering: the deep
3.7\,cm top margin exists because Crowdmark stamps a QR code into the top
band of \emph{every} page at upload -- interior pages keep the band empty,
and the cover lifts its NAME/ID box up into it, beside the QR.
